\documentclass[11pt]{article}
\usepackage[T1]{fontenc}
\usepackage[utf8]{inputenc}
\usepackage{lmodern}
\usepackage[margin=1in]{geometry}
\usepackage{amsmath,amssymb,amsthm,booktabs,microtype}
\usepackage{url}
\usepackage[colorlinks=true,linkcolor=blue,citecolor=blue,urlcolor=blue]{hyperref}
\hypersetup{pdftitle={The NNE distribution on Fuss--Catalan paths with a fixed number of peaks},
  pdfauthor={Xiaojian Zeng},pdfsubject={Fuss--Catalan enumeration and real-rootedness}}
\newtheorem{theorem}{Theorem}
\newtheorem{lemma}[theorem]{Lemma}
\DeclareMathOperator{\pk}{pk}
\DeclareMathOperator{\nne}{nne}
\newcommand{\N}{\mathsf{N}}
\newcommand{\E}{\mathsf{E}}
\newcommand{\W}{W^{(m)}_{n,k}}
\title{The NNE distribution on Fuss--Catalan paths with a fixed number of peaks}
\author{Xiaojian Zeng\\[3pt]
\small PBC School of Finance, Tsinghua University, 100190\\
\small \href{mailto:xiaojianzeng@foxmail.com}{xiaojianzeng@foxmail.com}}
\date{}
\begin{document}
\maketitle
\begin{abstract}
For paths from $(0,0)$ to $(n,mn)$ staying above $y=mx$, with $m\geq2$,
we enumerate consecutive NNE factors after fixing the number $k$ of NE
peaks. A first-return decomposition and Lagrange inversion give
\[
 W^{(m)}_{n,k}(t)=\frac{1}{n}\binom nk
       \sum_{j=1}^{k}\binom kj\binom{mn-k}{j-1}t^j.
\]
This is a positive scalar multiple of the ordinary Dyck-path polynomial
of semilength $mn+1$ with $k$ peaks, and hence proves the corresponding
Fuss--Catalan real-rootedness conjecture. A self-contained Jacobi-polynomial
argument shows that zero is simple and all other zeros are distinct and
strictly negative.
\end{abstract}

\section{Statement and relation to earlier work}

Fix integers $m\geq2$ and $n\geq1$. Let $\mathcal D^{(m)}_n$ be the
set of words containing $mn$ north steps $\N=(0,1)$ and $n$ east steps
$\E=(1,0)$, such that every prefix satisfies
$\#\N\geq m\#\E$. These are the paths from $(0,0)$ to $(n,mn)$
that stay in $y\geq mx$. For $D\in\mathcal D^{(m)}_n$, let $\pk(D)$
and $\nne(D)$ count the consecutive factors $\N\E$ and $\N\N\E$,
respectively. These counts do not wrap around the ends of the word.
Equivalently, each maximal north run contributes one peak, and contributes
one NNE factor exactly when its length is at least two. For $1\leq k\leq n$,
put
\begin{equation}\label{eq:def}
 \W(t)=\sum_{\substack{D\in\mathcal D^{(m)}_n\\\pk(D)=k}}t^{\nne(D)}.
\end{equation}

The ordinary Dyck-path joint distribution of peaks and NNE factors was
enumerated by B\'ona, Dimitrov, Labelle, Li, Pappe, Vindas-Mel\'endez and
Zhuang \cite[Theorem~2]{BonaEtAl2025}; their Theorem~27 proves
real-rootedness. Wang and Zhang \cite{WangZhang2024} reproduce these
results and establish the associated interlacing conjectures.
Alexandersson's catalog \cite{AlexanderssonCatalan} asks for the
real-rootedness extension to Fuss--Catalan paths in an $n\times nm$
rectangle. We resolve this question with the orientation and literal
factor conventions in \eqref{eq:def}.

\begin{theorem}\label{thm:main}
For every integer $m\geq2$, $n\geq1$ and $1\leq k\leq n$,
\begin{equation}\label{eq:formula}
 \W(t)=\frac{1}{n}\binom nk
       \sum_{j=1}^{k}\binom kj\binom{mn-k}{j-1}t^j.
\end{equation}
The polynomial has degree $k$, a simple zero at $0$, and $k-1$ distinct
strictly negative zeros. In particular, all its zeros are real.
\end{theorem}

Write $V_{r,k}(t)$ for the corresponding polynomial on ordinary Dyck
paths of semilength $r$, with the same NE and NNE conventions. The
enumeration in \eqref{eq:formula} gives the exact identity
\begin{equation}\label{eq:ordinary}
 \W(t)=\frac{\binom{n-1}{k-1}}{\binom{mn+1}{k-1}}\,V_{mn+1,k}(t).
\end{equation}
The contribution is this joint enumeration on the Fuss family and its
identification with the ordinary family. Generalized first-return
decompositions are standard; a nearby application is Zhang and Zhuang's
enumeration of augmented paths \cite{ZhangZhuang2024}. Their Lemma~12
marks long ascents after augmentation, whereas \eqref{eq:def} marks both
peaks and NNE factors before augmentation. The ordinary multi-pattern
refinement in \cite[Theorem~5]{BonaEtAl2025} also retains the ordinary
step set. No equivalent Fuss identification was found in these inspected
works; this comparison makes no exhaustive priority claim.

\section{The joint enumeration}

Give $\N$ height increment $1$ and $\E$ increment $-m$. A valid word
is then an excursion that stays at nonnegative height. We include the
empty excursion when forming generating functions.

\begin{lemma}\label{lem:split}
Every nonempty excursion has a unique decomposition
\begin{equation}\label{eq:split}
 D=\N P_1\N P_2\cdots\N P_m\E P_{m+1},
\end{equation}
where each $P_i$, translated to start at height zero, is an excursion,
possibly empty. The displayed $\E$ is the first return to height zero.
\end{lemma}

\begin{proof}
The first return must be an east step from height $m$. Before it, mark
the last north step from each height $i$ to $i+1$, for $0\leq i<m$.
Each exists because upward steps have size one. The marked steps occur
in increasing order of $i$: after the last departure from $i$, the path
still has to reach $m$ and therefore must later depart from every larger
level below $m$.

For $1\leq i<m$, let $P_i$ be the segment between the marked departures
from $i-1$ and $i$; let $P_m$ end just before the distinguished east step.
Each starts and ends at height $i$. If it visited a height below $i$,
its return to $i$ would require another north departure from $i-1$,
contradicting the choice of the marked step. This argument also applies
when an east step crosses several levels. The remaining suffix $P_{m+1}$
is an excursion at zero.

Conversely, inserting any $m+1$ excursions into \eqref{eq:split} yields
a valid word. Before the displayed east step its height is positive,
so that step is the first return. No subsequent component before this
return descends below its displayed base level. Thus the displayed
north steps are exactly the marked last departures, proving uniqueness
and the inverse construction.
\end{proof}

Every nonempty excursion starts with at least $m$ north steps and ends
with an east step. Since $m\geq2$, its first east step already has both
preceding letters inside the excursion. Consequently every NE or NNE
factor ending at an east step inside a component of \eqref{eq:split}
is counted internally. All additional factors end at the displayed east
step. The possibilities are exhaustive:
\begin{center}
\begin{tabular}{@{}lcc@{}}
\toprule
Component condition & Extra peaks & Extra NNE factors\\
\midrule
$P_m\ne\varnothing$ & $0$ & $0$\\
$P_m=\varnothing$, $P_{m-1}\ne\varnothing$ & $1$ & $0$\\
$P_m=P_{m-1}=\varnothing$ & $1$ & $1$\\
\bottomrule
\end{tabular}
\end{center}
Indeed, in the first case the displayed east step follows an east step;
in the second its two preceding letters are $\E\N$; in the third they
are $\N\N$. The last case includes $m=2$.

Let $z$ mark east steps, $u$ mark peaks and $t$ mark NNE factors, and set
\[
 F=1+\sum_{n\geq1}\sum_{D\in\mathcal D^{(m)}_n}
            z^n u^{\pk(D)}t^{\nne(D)}.
\]
The three rows of the table contribute, respectively,
$zF^m(F-1)$, $zuF^{m-1}(F-1)$ and $zutF^{m-1}$. Hence
\begin{equation}\label{eq:F}
 F=1+zF^{m-1}\bigl(F(F-1)+u(F-1)+ut\bigr).
\end{equation}
Set $Y=F-1$. Then $Y=z\phi(Y)$, where
\[
 \phi(v)=(1+v)^{m-1}\bigl(v(1+v)+u(v+t)\bigr).
\]
Formal Lagrange inversion \cite[Section~2.1]{Gessel2016} gives
\[
 [z^n]Y=\frac1n[v^{n-1}]\phi(v)^n.
\]
For example, this identity may be applied over $\mathbb Q(u,t)$, where
$\phi(0)=ut$ is invertible. The defining recursion has a unique
zero-constant-term solution whose coefficients belong to
$\mathbb Q[u,t]$, so subsequent coefficient extraction is legitimate;
no analytic convergence is involved.

Selecting $u^k$ chooses $k$ of the $n$ factors $u(v+t)$ and yields
\begin{align}
 [z^nu^k]Y
 &=\frac1n\binom nk[v^{n-1}]
       (1+v)^{(m-1)n}v^{n-k}(1+v)^{n-k}(v+t)^k\notag\\
 &=\frac1n\binom nk[v^{k-1}](1+v)^{mn-k}(v+t)^k.\label{eq:extraction}
\end{align}
The coefficient of $t^j$ in the last factor is $\binom kj v^{k-j}$.
The remaining exponent of $v$ is $j-1$; the $j=0$ term vanishes.
This proves \eqref{eq:formula}. Since $mn-k\geq k$, every displayed
coefficient is positive. In particular,
\[
 [t]\W(t)=\binom{n-1}{k-1}>0,
\]
and the degree is $k$. For $k=1$ the formula gives $\W(t)=t$.
As a consistency check, Vandermonde's identity gives
\[
 \W(1)=\frac1n\binom nk\binom{mn}{k-1}.
\]

\section{Real zeros and simplicity}

The ordinary formula of B\'ona et al., as stated in
\cite[Theorem~1]{WangZhang2024}, implies, for $r\geq2k+1$,
\[
 V_{r,k}(t)=\frac1k\binom r{k-1}
                \sum_{j=1}^k\binom kj\binom{r-k-1}{j-1}t^j.
\]
Here the ordinary up and down letters correspond directly to $\N$ and
$\E$, with no reversal of factors. Substituting $r=mn+1\geq2k+1$
and using $\frac{k}{n}\binom nk=\binom{n-1}{k-1}$ proves
\eqref{eq:ordinary}. The scalar is positive, so ordinary real-rootedness
\cite[Theorem~27]{BonaEtAl2025} already proves real-rootedness in
Theorem~\ref{thm:main}.

We give a second, self-contained root argument to establish simplicity
as well. Put
\[
 d=k-1,\qquad A=mn-k,\qquad b=A-d=mn-2k+1\geq1.
\]
Define
\begin{equation}\label{eq:J}
 J(x)=\sum_{\ell=0}^d\binom{d+1}{d-\ell}\binom{d+b}{\ell}
           \left(\frac{x-1}{2}\right)^\ell
           \left(\frac{x+1}{2}\right)^{d-\ell}.
\end{equation}
This is the classical Jacobi polynomial $P_d^{(1,b)}$;
see \cite[Eq.~18.5.8 and Table~18.5.1]{DLMFJacobi} for the normalization
and Rodrigues formula. Its leading coefficient is positive, so $\deg J=d$.
The product rule gives
\begin{equation}\label{eq:Rodrigues}
 (1-x)(1+x)^bJ(x)=\frac{(-1)^d}{2^d d!}
       \frac{d^d}{dx^d}\bigl((1-x)^{d+1}(1+x)^{d+b}\bigr).
\end{equation}
To check the identity directly, differentiate the first factor $d-\ell$
times and the second factor $\ell$ times. After division by the weight,
the resulting summand is
\[
 2^{-d}(-1)^\ell\binom{d+1}{d-\ell}\binom{d+b}{\ell}
                  (1-x)^\ell(1+x)^{d-\ell},
\]
which is the corresponding summand in \eqref{eq:J}.

Let $q$ have degree less than $d$. Multiply \eqref{eq:Rodrigues} by $q$
and integrate over $[-1,1]$. Integrating by parts $d$ times gives
\begin{equation}\label{eq:orthog}
 \int_{-1}^1 q(x)J(x)(1-x)(1+x)^b\,dx=0.
\end{equation}
All boundary terms vanish because
$(1-x)^{d+1}(1+x)^{d+b}$ has zeros of order at least $d+1$ at both
endpoints; the remaining integral contains $q^{(d)}=0$.

Suppose $J$ had only $s<d$ sign-changing roots $x_1,\ldots,x_s$ in
$(-1,1)$. Taking $q(x)=\prod_{i=1}^s(x-x_i)$ makes $qJ$ have one fixed,
nonzero sign away from its finitely many zeros. Since the weight in
\eqref{eq:orthog} is strictly positive in the interval, this contradicts
\eqref{eq:orthog}. Thus $J$ has $d$ distinct roots, all in $(-1,1)$.
For $d=0$, $J=1$ and the assertion is vacuous.

Termwise substitution into \eqref{eq:J} now gives
\begin{equation}\label{eq:transform}
 (1-t)^d J\!\left(\frac{1+t}{1-t}\right)
       =\sum_{\ell=0}^d\binom{k}{\ell+1}\binom A\ell t^\ell.
\end{equation}
All denominators cancel, so this is a polynomial identity. Each zero
$x_i\in(-1,1)$ of $J$ gives the distinct negative zero
$t_i=(x_i-1)/(x_i+1)$ on the right. These $d$ zeros exhaust its degree.
Its constant term is $k$, and its value at $t=1$ is positive.
Multiplying \eqref{eq:transform} by $\frac1n\binom nk t$ proves the
remaining assertions of Theorem~\ref{thm:main}.

The lower parameter boundary $m=2$, $k=n$ gives $b=1$, so the same
positive-weight proof applies there. The functional equation
\eqref{eq:F} is specific to $m\geq2$: an ordinary component may start
with NE, allowing an NNE factor to cross its beginning. The ordinary
result is used through \eqref{eq:ordinary}, rather than by specializing
\eqref{eq:F} to $m=1$.

\section{Exact checks and reproducibility}

The proof above is independent of computation. The accompanying exact
verifier checks the factor conventions, coefficients and root counts.
Its implementation appends one north or east letter at a time and
retains the numbers of each letter, the last two letters, and the two
statistics. It checked all $1\leq k\leq n$ for
\[
 m\in\{2,3,4,7\},\quad 1\leq n\leq10,
\]
and for the nine additional pairs
\[
\begin{gathered}
 (2,14),(2,16),(3,13),(5,8),(6,11),\\
 (11,8),(17,5),(31,5),(100,3).
\end{gathered}
\]
This gives $49$ pairs and $303$ polynomials. All coefficients agree with
\eqref{eq:formula}. For $Q(t)=\W(t)/t$, exact signed Euclidean remainder
chains count $k-1$ roots in $(-\infty,0)$ by Sturm's theorem
\cite[Theorem~2.50]{BasuPollackRoy2006}, and the terminal gcd has degree
zero. The latter check has the derivative Bézout interpretation
$aQ+cQ'=1$ for suitable rational polynomials $a,c$, when $Q$ is
nonconstant. The corresponding normalized derivative Bézout pair is
described by Zeng \cite[Eq.~(1) and Lemma~2.1]{Zeng2026Bezout}.
Only the coprimality test is needed here, and it is performed
by exact polynomial division. For $k=1$, $Q=1$ is handled directly.

The accompanying \texttt{verify.py} reproduces the $303$-polynomial
check using only the Python standard library and additionally verifies
eleven instances of \eqref{eq:Rodrigues} and \eqref{eq:transform} after
clearing denominators. It records all coefficients, endpoint signs and
root counts in \texttt{verification.json}; run it with
\texttt{python3 verify.py}. The eleven $(d,b)$ inputs for the identity
checks are explicitly listed in the script.
These are exact finite observations; Theorem~\ref{thm:main} follows from
the general decomposition, coefficient extraction and orthogonality
argument.

\paragraph{Computational and AI assistance.}
OpenAI Codex assisted with proof exploration, verification programming,
literature checks and manuscript drafting. SymPy~1.14.0
\cite{MeurerEtAl2017} was used for symbolic differentiation and expansion
during proof exploration. A separate model pass examined
the counting and root arguments and implemented the letter-based verifier.
This assistance is not external human peer review. No proof-assistant
formalization or floating-point root fit is used in the argument.

\bibliographystyle{plainurl}
\bibliography{references}
\end{document}
